\par\Needspace{6\baselineskip}\medskip
\hypertarget{native-081-article062}{}\pdfbookmark[2]{62.}{article-062}
\begin{center}62.\end{center}
Die heliozentrische Länge des Himmelskörpers sei $\lambda$ und seine Breite $\beta$, seine geozentrische Länge $l$ und seine Breite $b$, sein Abstand von der Sonne $r$ und von der Erde $\Delta$. Die heliozentrische Länge der Erde sei schließlich $L$, ihre Breite $B$ und ihr Abstand von der Sonne $R$. Da wir $B=0$ nicht voraussetzen, gelten unsere Formeln auch, wenn man die heliozentrischen Orte und den geozentrischen Ort auf eine beliebige andere Ebene als die Ekliptik bezieht. Man muss lediglich die Bezeichnungen Länge und Breite weglassen. Außerdem lässt sich die Parallaxe sofort berücksichtigen. Dazu genügt es, den heliozentrischen Ort der Erde unmittelbar auf einen Punkt ihrer Oberfläche anstelle ihres Mittelpunkts zu beziehen. Wir setzen ferner $r\cos\beta=r'$, $\Delta\cos b=\Delta'$, $R\cos B=R'$. Bezieht man nun die Orte des Himmelskörpers und der Erde im Raum auf drei Ebenen, von denen die erste die Ekliptik ist und die zweite und dritte ihre Pole auf den Längen $N$ und $N+90^0$ haben, so ergeben sich unmittelbar die folgenden Gleichungen
\[
\begin{gathered}
r'\cos(\lambda-N)-R'\cos(L-N)=\Delta'\cos(l-N)\\
r'\sin(\lambda-N)-R'\sin(L-N)=\Delta'\sin(l-N)\\
r'\operatorname{tang}\beta-R'\operatorname{tang}B=\Delta'\operatorname{tang}b,
\end{gathered}
\]
wobei der Winkel $N$ völlig beliebig ist. Die erste und zweite Gleichung liefern sofort $l-N$ und $\Delta'$. Aus diesen Größen ergibt sich mit der dritten Gleichung $b$, und aus $b$ und $\Delta'$ erhält man $\Delta$. Um die Rechnung möglichst bequem zu gestalten, wählen wir den beliebigen Winkel $N$ auf die folgenden drei Arten.

I. Mit $N=L$ setzen wir $\frac{r'}{R'}\sin(\lambda-L)=P$ und $\frac{r'}{R'}\cos(\lambda-L)-1=Q$. Dann findet man $l-L$, $\frac{\Delta'}{R'}$ und $b$ mit den Formeln
\[
\begin{gathered}
\operatorname{tang}(l-L)=\frac{P}{Q}\\[.5em]
\frac{\Delta'}{R'}=\frac{P}{\sin(l-L)}=\frac{Q}{\cos(l-L)}\\[.5em]
\operatorname{tang}b=\frac{\frac{r'}{R'}\operatorname{tang}\beta-\operatorname{tang}B}{\frac{\Delta'}{R'}}\mathord{\cdot}
\end{gathered}
\]

\SourceLeaf{082}II. Mit $N=\lambda$ setzen wir $\frac{R'}{r'}\sin(\lambda-L)=P$ und $1-\frac{R'}{r'}\cos(\lambda-L)=Q$. Dann gilt
\[
\begin{gathered}
\operatorname{tang}(l-\lambda)=\frac{P}{Q}\\[.5em]
\frac{\Delta'}{r'}=\frac{P}{\sin(l-\lambda)}=\frac{Q}{\cos(l-\lambda)}\\[.5em]
\operatorname{tang}b=\frac{\operatorname{tang}\beta-\frac{R'}{r'}\operatorname{tang}B}{\frac{\Delta'}{r'}}\mathord{\cdot}
\end{gathered}
\]

III. Mit $N=\frac12(\lambda+L)$ findet man $l$ und $\Delta'$ mit den Gleichungen
\[
\begin{gathered}
\operatorname{tang}\{l-\tfrac12(\lambda+L)\}=\frac{r'+R'}{r'-R'}\operatorname{tang}\tfrac12(\lambda-L)\\[.5em]
\Delta'=\frac{(r'+R')\sin\frac12(\lambda-L)}{\sin(l-\frac12(\lambda+L))}=\frac{(r'-R')\cos\frac12(\lambda-L)}{\cos(l-\frac12(\lambda+L))}
\end{gathered}
\]
und anschließend $b$ mit der oben angegebenen Gleichung. Der Logarithmus des Bruchs $\frac{r'+R'}{r'-R'}$ lässt sich bequem berechnen, wenn man $\frac{R'}{r'}=\operatorname{tang}\zeta$ setzt. Daraus folgt $\frac{r'+R'}{r'-R'}=\operatorname{tang}(45^0+\zeta)$. Auf diese Weise ist Verfahren III zur Bestimmung von $l$ etwas kürzer als I und II. Für die übrigen Rechenschritte halten wir jedoch I und II für zweckmäßiger als III.
